\documentclass[11pt,letterpaper]{article}
% preamble.tex — Shared preamble for all documents
\usepackage{graphicx}
\usepackage{hyperref}
\usepackage{enumitem}
\usepackage{titlesec}
\usepackage{fancyhdr}
\usepackage{setspace}
\usepackage{parskip}
\usepackage{etoolbox}
\usepackage{xspace}
\usepackage{needspace}
\usepackage{changepage}
\usepackage{geometry}
\usepackage{ragged2e}
\usepackage{environ}

% Hyperref setup
\hypersetup{
  colorlinks=false,
  pdfborder={0 0 0},
  pdftitle={The Letter to the Father},
  pdfauthor={},
}

% Load project files
% colors.tex — Color palette for The Letter to the Father
\usepackage{xcolor}

% Primary palette — muted, period-appropriate, prestige
\definecolor{kafkablack}{HTML}{1A1A1A}
\definecolor{kafkagray}{HTML}{4A4A4A}
\definecolor{kafkalight}{HTML}{8A8A8A}
\definecolor{kafkacream}{HTML}{F5F0E8}
\definecolor{kafkasepia}{HTML}{704214}
\definecolor{kafkablood}{HTML}{6B1C23}
\definecolor{kafkaink}{HTML}{2C2C3A}

% Accent colors — for package materials
\definecolor{prestigegold}{HTML}{B8860B}
\definecolor{sectionrule}{HTML}{3A3A3A}
\definecolor{pullquote}{HTML}{5A4A3A}

% metadata.tex — Project metadata for The Letter to the Father
\newcommand{\filmtitle}{The Letter to the Father}
\newcommand{\filmtitleupper}{THE LETTER TO THE FATHER}
\newcommand{\filmsubtitle}{An Original Screenplay}
\newcommand{\filmauthor}{Based on the life and writings of Franz Kafka}
\newcommand{\filmyear}{2026}
\newcommand{\filmdraft}{Revised Draft -- Public Development Example}
\newcommand{\filmdate}{October 2026 -- Example Edition}
\newcommand{\filmlogline}{In the final years of his life, a frail and tormented Franz Kafka struggles to write, love, and exist within a world of family pressure, bureaucratic humiliation, illness, and spiritual dread --- while the stories he creates begin to mirror and consume the life he can no longer control.}
\newcommand{\filmgenre}{Prestige Drama / Literary Biopic}
\newcommand{\filmrating}{R}
\newcommand{\filmlead}{Adrien Brody}
\newcommand{\filmbudgetrange}{\$12--18M}
\newcommand{\filmruntime}{Approximately 110 minutes}


% Public example status travels with every downloadable document.
\AddToHook{shipout/foreground}{\put(40,-776){\makebox(0,0)[l]{\fontsize{8}{9}\selectfont\textit{Development example. All casting, financing and awards language is illustrative; no attachments confirmed.}}}}

% typography.tex — Typography settings
\usepackage{fontspec}
\usepackage{microtype}

% Prefer repo-local metric-compatible fonts so Tectonic does not depend on
% absolute host-system font paths. Fall back to system fonts only if the vendored
% font payload is unavailable.
\IfFileExists{../../../fonts/Tinos-Regular.ttf}{
  \setmainfont{Tinos-Regular.ttf}[
    Path=../../../fonts/,
    BoldFont=Tinos-Bold.ttf,
    ItalicFont=Tinos-Italic.ttf,
    BoldItalicFont=Tinos-BoldItalic.ttf,
  ]
}{
  \IfFontExistsTF{TeX Gyre Termes}{
    \def\projectmainfont{TeX Gyre Termes}
  }{
    \IfFontExistsTF{Times New Roman}{
      \def\projectmainfont{Times New Roman}
    }{
      \def\projectmainfont{Latin Modern Roman}
    }
  }
  \setmainfont{\projectmainfont}
}

\IfFileExists{../../../fonts/Arimo-Regular.ttf}{
  \setsansfont{Arimo-Regular.ttf}[
    Path=../../../fonts/,
    BoldFont=Arimo-Bold.ttf,
    ItalicFont=Arimo-Italic.ttf,
    BoldItalicFont=Arimo-BoldItalic.ttf,
  ]
}{
  \IfFontExistsTF{TeX Gyre Heros}{
    \def\projectsansfont{TeX Gyre Heros}
  }{
    \IfFontExistsTF{Helvetica}{
      \def\projectsansfont{Helvetica}
    }{
      \def\projectsansfont{Latin Modern Sans}
    }
  }
  \setsansfont{\projectsansfont}
}

\IfFileExists{../../../fonts/Cousine-Regular.ttf}{
  \setmonofont{Cousine-Regular.ttf}[
    Path=../../../fonts/,
    BoldFont=Cousine-Bold.ttf,
    ItalicFont=Cousine-Italic.ttf,
    BoldItalicFont=Cousine-BoldItalic.ttf,
    Scale=0.95,
  ]
  \def\screenplayfont{%
    \fontspec{Cousine-Regular.ttf}[%
      Path=../../../fonts/,
      BoldFont=Cousine-Bold.ttf,
      ItalicFont=Cousine-Italic.ttf,
      BoldItalicFont=Cousine-BoldItalic.ttf,
      Scale=1.0,
    ]%
  }
}{
  \IfFontExistsTF{Courier New}{
    \def\projectmonofont{Courier New}
  }{
    \IfFontExistsTF{Courier}{
      \def\projectmonofont{Courier}
    }{
      \def\projectmonofont{Latin Modern Mono}
    }
  }
  \setmonofont{\projectmonofont}[Scale=0.95]
  \def\screenplayfont{\fontspec{\projectmonofont}[Scale=1.0]}
}

% Microtype settings
\microtypesetup{
  protrusion=true,
  expansion=false,
}

% studio_package_macros.tex — Formatting macros for the studio package

% ============================================================
% SECTION STYLING
% ============================================================
\titleformat{\section}
  {\Large\bfseries\color{kafkablack}\sffamily}
  {}
  {0em}
  {}
  [\vspace{2pt}\textcolor{sectionrule}{\rule{\textwidth}{0.5pt}}\vspace{6pt}]

\titleformat{\subsection}
  {\large\bfseries\color{kafkagray}\sffamily}
  {}
  {0em}
  {}

\titleformat{\subsubsection}
  {\normalsize\bfseries\color{kafkagray}\sffamily}
  {}
  {0em}
  {}

% ============================================================
% PULL QUOTE
% ============================================================
% Usage: \pullquote{Text here}
\NewEnviron{pullquote}{%
  \vspace{8pt}%
  \begin{adjustwidth}{0.5in}{0.5in}%
    \color{pullquote}\itshape\large%
    \BODY%
  \end{adjustwidth}%
  \vspace{8pt}%
}

% ============================================================
% HIGHLIGHT BOX
% ============================================================
\newcommand{\highlightbox}[1]{%
  \par\vspace{6pt}%
  \noindent\colorbox{kafkacream}{%
    \parbox{\dimexpr\textwidth-2\fboxsep}{%
      \vspace{4pt}#1\vspace{4pt}%
    }%
  }%
  \par\vspace{6pt}%
}

% ============================================================
% TARGET ENTRY (for director/producer lists)
% ============================================================
% Usage: \target{Name}{Why they fit}{Tier}
\newcommand{\target}[3]{%
  \needspace{4\baselineskip}%
  \subsubsection*{#1 \hfill \textnormal{\small\textit{#3}}}%
  #2\par\vspace{6pt}%
}

% ============================================================
% SCENE CARD (for showcase section)
% ============================================================
% Usage: \scenecard{Title}{Description}
\newcommand{\scenecard}[2]{%
  \needspace{5\baselineskip}%
  \vspace{4pt}%
  \noindent\textbf{\sffamily #1}\\[3pt]%
  #2\par\vspace{8pt}%
  \textcolor{sectionrule}{\rule{0.3\textwidth}{0.25pt}}%
  \vspace{8pt}%
}

% ============================================================
% TAGLINE
% ============================================================
\newcommand{\tagline}[1]{%
  \vspace{4pt}%
  \begin{center}%
    {\large\itshape\color{kafkagray} ``#1''}%
  \end{center}%
  \vspace{4pt}%
}

% ============================================================
% COMP FILM ENTRY
% ============================================================
% Usage: \compfilm{Title (Year)}{Relevance}
\newcommand{\compfilm}[2]{%
  \noindent\textbf{#1}\\%
  {\small #2}\par\vspace{6pt}%
}

% ============================================================
% RISK ENTRY
% ============================================================
% Usage: \risk{Risk}{Mitigation}
\newcommand{\risk}[2]{%
  \noindent\textbf{Risk:} #1\\%
  \textbf{Mitigation:} #2\par\vspace{8pt}%
}

% ============================================================
% PACKAGE GEOMETRY
% ============================================================
\newcommand{\packagegeometry}{%
  \newgeometry{
    left=1in,
    right=1in,
    top=1in,
    bottom=1in,
  }%
}

\geometry{left=0.75in, right=0.75in, top=0.75in, bottom=0.75in}
\pagestyle{empty}
\setlength{\parindent}{0pt}
\setlength{\parskip}{5pt}

\begin{document}

\begin{center}
{\Large\bfseries\sffamily \filmtitleupper}\\[3pt]
{\normalsize\sffamily\color{kafkagray} Awards Campaign Brief}
\end{center}
\vspace{2pt}
\textcolor{sectionrule}{\rule{\textwidth}{0.5pt}}
\vspace{4pt}

\textbf{The Primary Case}\\
This is a Best Actor film first. Everything else begins there. Adrien Brody as Kafka is a role with the exact combination awards campaigns need: recognizable historical subject, severe physical and emotional demands, multiple clip-ready scenes, and a career narrative simple enough to explain in one sentence. The campaign writes itself: \textit{The Pianist}, the long wilderness, \textit{The Brutalist}, then the definitive performance of the late phase.

\vspace{2pt}

\textbf{Why The Brody Narrative Is So Strong}\\
The role is built around contradiction, which is what makes it campaignable. Kafka can make an institution tell the truth but cannot make his own father read a letter. Brody has to play public intelligence, private shame, social terror, romantic tenderness, philosophical combat, and terminal physical decline in the same film. That is unusually broad evidence for a Best Actor push, and it arrives with the added advantage of a role that feels culturally inevitable once named.

\vspace{2pt}

\textbf{The Clip Set}\\
\textit{The Elevator}: a bureaucratic nothing turning into a metaphysical event, carried almost entirely in the face. --- \textit{The Engagement Dinner}: one of the year's best social-collapse scenes. --- \textit{The Letter Returns}: the title object comes back undelivered and the whole movie tightens around it. --- \textit{Burn Everything}: a friendship scene, moral argument, and self-sentencing at once. --- \textit{The Final Notes}: physical devastation, humor, and spiritual clarity in a man who cannot speak.

\vspace{2pt}

\textbf{Secondary Categories}\\
Best Picture becomes plausible if the film is campaigned as a pressure-system drama rather than a generic literary biopic. Best Director is attachment-dependent but immediately real with the right filmmaker. Best Adapted Screenplay is strategically cleaner than Original if the final crediting uses Kafka's writings and letters. Cinematography, Production Design, and Supporting Actor for Hermann are all credible adjacent lanes, which matters because campaigns land better when the film has more than one point of entry.

\vspace{2pt}

\textbf{How to Talk About It}\\
Not ``important'' or ``literary.'' Those are dead words. The effective language is: career-best, devastating restraint, formally controlled but emotionally violent, a biographical drama with the pressure of a thriller, the rare literary portrait with real propulsion. The office line, father-letter engine, and late Berlin warmth are not just story virtues; they are campaign language because they help explain why the movie moves rather than merely broods.

\vspace{2pt}

\textbf{Why It Can Survive the Risks}\\
``Another literary biopic'' fatigue is real, but this is structurally cleaner and more marketable than that label implies. ``Too European'' is mitigated by Brody, family pressure, illness, and the universality of being unread by the people closest to you. A crowded Best Actor field is always possible, but the role's breadth, the clip variety, and the career-culmination narrative give the film multiple ways to stay alive in the conversation.

\vspace{2pt}

\textbf{Ideal Campaign Shape}\\
Venice launch, then Telluride/TIFF acceleration, then a disciplined fall release and FYC corridor. The movie wants curation first, then amplification. If the festival response lands where it should, the campaign story is unusually easy to sustain because the role, the title, and the central contradiction are all memorable on first contact.

\vspace{4pt}
\textcolor{sectionrule}{\rule{\textwidth}{0.5pt}}
\vspace{2pt}
\begin{center}
{\small\color{kafkalight}\textit{An awards play with one undeniable center, five clean clip lanes, and a campaign narrative that does not need to be invented.}}
\end{center}

\end{document}
